\section{Smoothing from the actual Hilbert quotient to Meyer's quotient}
\label{sec:hilbert-meyer-bridge}

The comparison in this section is an operator factorization.  It does
not identify two traces by first assigning both the same zero sum.
It gives a nuclear smoothing map from the specified Hilbert quotient
to the actual smooth quotient of MT-M24, and then uses trace cyclicity.
The uniform-interior hypothesis of MT-M26 is retained throughout.
Consequently the comparison is unconditional on a wide strip and
applies to every positive narrow Dirichlet strip using $A$.  It is
not an independent proof that all zeros belong to those narrow strips.

\subsection{The spaces and the one hypothesis about the zero divisor}

Fix one primitive unitary Hecke function and the completed nontrivial
zero divisor of $\Xi_\eta$.  Let $\mathcal E$ and its closed ideal
$\mathcal I_\eta$ be those of GRH\_NIGHT MT-M18, and put
$Q=\mathcal E/\mathcal I_\eta$, with quotient map $q$.
For $a\in\mathbb R$, $s\ge0$ and $I=[u,v]$, $u<v$, use the actual
spaces of MT-M26:
\[
 \begin{aligned}
 H&=L^2\!\left(\mathbb R,
     \max(e^{2ux},e^{2vx})(1+x^2)^{s/2}\,dx\right),\\
 J&=\overline{T_a\mathcal I_\eta}^{\,H},
       &T_af(x)&=e^{ax}f(x),\\
 Z&=H/J,\qquad
       i_a:Q\longrightarrow Z,\quad [f]\longmapsto[T_af].
 \end{aligned}
 \tag{MA-HM.1}
\]
The hypothesis is that a closed strip contains the real parts of
every nontrivial completed zero and has a positive margin inside
the available Mellin strip:
\[
       \operatorname{Re}\rho\in[c_0,c_1]\quad\hbox{for every }\rho,
       \qquad [c_0-a,c_1-a]\subset(u,v).
 \tag{MA-HM.2}
\]
This is exactly MT-M26 (M26.2), not a condition on only the zeros
already selected inside $I$.  The closed interval $[0,1]$ works
unconditionally.  At $a=1/2$ and $I=[-\alpha,\alpha]$, it works
unconditionally for $\alpha>1/2$, and for every $\alpha>0$ under
the RH of the particular function.  In the Dirichlet case that
last input is $A$; the general sufficient input is $M$.

Choose the adapted rectangles $\Omega_n$ strictly between the
edges $a+u,a+v$ as in the proof of MT-M26, and put $\nu_n=2+|n|$.
The exact inherited data used below are
\[
 \begin{aligned}
 d_n&=\dim J_n=O(\log\nu_n),\\
 R_m\mathcal T_n&=\delta_{mn}\operatorname{id}_{J_n},
       &\mathcal T_n&:J_n\longrightarrow\mathcal E,\\
 p_k(\mathcal T_n z)&\le C_k\nu_n^{B_k}\|z\|_n,\\
 r_n&:Z\longrightarrow J_n,\qquad
       \sup_n\|r_n\|\le C,\qquad r_ni_a=R_n .
 \end{aligned}
 \tag{MA-HM.3}
\]
Here $p_k$ is any enumeration of the defining exponential Schwartz
seminorms.  The $J_n$ are the quotient holomorphic supremum-norm
algebras with their actual full zero jets; $R_n$ is the block map
on $Q$.  MT-M18 identifies $Q$ topologically with the rapid block
space $s(J)$ and supplies the lifts.  MT-M26 (M26.5)--(M26.9)
supplies the same data in the narrower adapted rectangles and
the uniform Hilbert restriction bound.  In particular no
unstable basis of individual zero evaluations is being substituted
for these stable finite algebras.

For $g\in\mathcal E$, let $M_{g,n}$ be multiplication in $J_n$
by the class of $w\mapsto\mathcal Mg(w-a)$.  Set
\[
 b_n(g)=\sup_{w\in\partial\Omega_n}|\mathcal Mg(w-a)|.
 \qquad
 \|M_{g,n}\|\le b_n(g),\qquad
 b_n(g)\le C_A(g)\nu_n^{-A}\quad(A>0).
 \tag{MA-HM.4}
\]
Every $C_A(g)$ can be bounded by finitely many seminorms of $g$,
and is bounded on bounded subsets of $\mathcal E$, by the
integration-by-parts estimate in MT-M26 (M26.12).

\subsection{A lift into the smooth test space, not just into its completion}

\paragraph{MA-HM-001 (order-zero nuclear smoothing through the actual quotient).}
Under (MA-HM.2), the series
\[
 \widetilde S_g z
       =\sum_{n\in\mathbb Z}\mathcal T_nM_{g,n}r_nz
       \quad(z\in Z),\qquad
 S_g=q\widetilde S_g:Z\longrightarrow Q
 \tag{MA-HM.5}
\]
converges absolutely in every seminorm of $\mathcal E$.
It defines a continuous linear map into $\mathcal E$, and is
of order zero: it admits a nuclear rank-one representation with
coefficients in $\ell^p$ for every $p>0$, bounded target vectors
in $\mathcal E$, and bounded functionals on $Z$.
The maps $g\mapsto\widetilde S_g$ are continuous for uniform
convergence on the unit ball of $Z$ in every seminorm of
$\mathcal E$.  Bounded test families give bounded order-zero
nuclear families.  These assertions also hold for $S_g$.

Writing $k_g(x)=e^{-ax}g(x)$, one has identities of actual operators
\[
             S_gi_a=C_{k_g}^{Q},\qquad
             i_aS_g=C_g^{Z}.
 \tag{MA-HM.6}
\]
Thus $C_g^Z$ factors through the stronger smooth quotient topology,
although the unsmoothed dense injection $i_a$ is not a topological
isomorphism onto $Z$.  The quotient-valued map $S_g$ is independent
of the chosen block lifts.  It is a genuine smoothed norm
intertwiner: with $U_yf(x)=f(x-y)$ on both radial models,
\[
             S_gU_y^Z=e^{-ay}U_y^QS_g .
 \tag{MA-HM.6a}
\]
Only the chosen representative lift $\widetilde S_g$ into
$\mathcal E$ need depend on the interpolation section.

\emph{Proof.}
For every seminorm $p_k$ and every $z\in Z$, (MA-HM.3)--(MA-HM.4)
give
\[
 \sum_n p_k(\mathcal T_nM_{g,n}r_nz)
      \le CC_k\|z\|_Z\sum_n\nu_n^{B_k}b_n(g)<\infty.
 \tag{MA-HM.7}
\]
The completeness of $\mathcal E$ gives a unique sum.  Choosing
$A>B_k+2$ in (MA-HM.4) bounds its $p_k$ seminorm by
a constant times $C_A(g)\|z\|_Z$.  This proves continuity and
the asserted continuous dependence on $g$.  It also justifies
applying every continuous block restriction to the sum.

Per-seminorm convergence alone would not be an adequate
substitute for a nuclear tensor representation, so we give that
representation explicitly.  Choose an Auerbach basis
$e_{n,j},e_{n,j}^*$ of each nonzero $J_n$,
$1\le j\le d_n$, with both norms equal to one.  Put
\[
     \ell_{n,j}=e_{n,j}^*r_n\in Z',\qquad
           \|\ell_{n,j}\|\le C.
 \tag{MA-HM.8}
\]
Then the finite-dimensional identity in each block yields
\[
       \widetilde S_g
           =\sum_{n,j}\mathcal T_nM_{g,n}e_{n,j}
                                      \otimes\ell_{n,j}.
 \tag{MA-HM.9}
\]
If $b_n(g)>0$, write its summand as
$\lambda_{n,j}v_{n,j}\otimes\ell_{n,j}$ with
\[
   \lambda_{n,j}=\sqrt{b_n(g)},\qquad
   v_{n,j}=b_n(g)^{-1/2}\mathcal T_nM_{g,n}e_{n,j}.
 \tag{MA-HM.10}
\]
If $b_n(g)=0$, the block multiplier is zero and its summands
are omitted.  For each fixed $k$,
\[
       p_k(v_{n,j})\le C_k\nu_n^{B_k}\sqrt{b_n(g)}.
 \tag{MA-HM.11}
\]
Thus the entire family of vectors is bounded in $\mathcal E$.
For every $p>0$,
\[
       \sum_{n,j}|\lambda_{n,j}|^p
          =\sum_nd_n b_n(g)^{p/2}<\infty.
 \tag{MA-HM.12}
\]
These are precisely the bounded-vector and equibounded-functional
conditions for an order-zero nuclear map.  If a Banach-disk
description is desired, take the closed absolutely convex hull
of the bounded target vectors: completeness of the Fr\'echet
space makes it a complete bounded disk.  Formula (MA-HM.9)
then factors through that Banach disk and the indicated
$\ell^p$ coefficients.  No common polynomial exponent for
all Schwartz seminorms is required.

For a bounded test set $\mathcal B\subset\mathcal E$, replace
$b_n(g)$ in the normalizing square root by
$b_n^{\mathcal B}=\sup_{g\in\mathcal B}b_n(g)$.
This is again a rapidly decreasing sequence.  The target
vectors in (MA-HM.10), now with $g$ varying over $\mathcal B$,
belong to one bounded disk by (MA-HM.11); the functionals
in (MA-HM.8) are fixed and uniformly bounded; and
(MA-HM.12) gives one summable coefficient sequence for each
$p>0$.  This proves the asserted boundedness in the
order-zero operator bornology.  Applying the continuous
quotient map $q$ proves the same assertions for $S_g$.

To prove the first operator identity in (MA-HM.6), let $[f]\in Q$.
By (MA-HM.3), the $n$th block of $S_gi_a[f]$ is
$M_{g,n}R_n[f]$.  On the other hand
$\mathcal M k_g(w)=\mathcal Mg(w-a)$, so the same is true
of $C_{k_g}^Q[f]$.  The block isomorphism $Q\simeq s(J)$
makes the two classes equal.  This proves an equality in
the actual quotient, not merely equality of a list of scalar
eigenvalues.

For the second identity, convolution is bounded on $H$ by
the translation estimate (M26.10).  It preserves $J$ since
\[
       C_gT_af=T_aC_{k_g}f
       \quad(f\in\mathcal E),
 \tag{MA-HM.13}
\]
and $\mathcal I_\eta$ is a convolution ideal.  Consequently
$C_g^Z$ is the bounded operator induced from the literal
ambient convolution.  The same equation proves
$C_g^Zi_a=i_aC_{k_g}^Q$.  Composing the first identity in
(MA-HM.6) with $i_a$ therefore shows that $i_aS_g$ and
$C_g^Z$ agree on $i_aQ$.  That image is dense, because
$T_a\mathcal E$ contains $C_c^\infty(\mathbb R)$, dense in
the weighted $L^2$ space.  Both operators on $Z$ are
bounded, so they agree on all of $Z$.

If two continuous maps $Z\to Q$ have the first identity in
(MA-HM.6), they agree on the dense subspace $i_aQ$, hence
everywhere because $Q$ is Hausdorff.  This proves the
independence of $S_g$.  The same argument proves (MA-HM.6a):
$U_y^Zi_a=e^{-ay}i_aU_y^Q$ by the definition of $T_a$,
and smooth convolution commutes with $U_y^Q$.  The two
sides of (MA-HM.6a) are continuous maps $Z\to Q$ and
therefore agree everywhere once they agree on $i_aQ$.

Only the available data (MA-HM.3) used (MA-HM.2).
No scalar trace formula, simplicity assumption, or positivity
was used in this proof.  The construction proves the
factorization even if one uses only the continuity and
density of $i_a$ before invoking the injectivity supplied
by the rest of MT-M26. \hfill$\square$

\subsection{Trace cyclicity and the actual Meyer common-space expression}

\paragraph{MA-HM-002 (the Hilbert trace is the actual Meyer quotient trace).}
Under the same hypothesis, one has
\[
 \operatorname{Tr}_{Z}C_g^Z
       =\operatorname{tr}_{Q}C_{k_g}^Q
       \qquad(g\in\mathcal E),
 \tag{MA-HM.14}
\]
where the left side is the Hilbert trace and the right side
is the canonical nuclear trace on the smooth quotient.
For $g\in C_c^\infty(\mathbb R)$ let $\beta$ be the idele
character in the inverse-Euler convention of MT-M22 and set
\[
       f_{g,a}(b)=\beta(b)^{-1}|b|^{-a}g(\log|b|)
                \quad(b\in C_K).
 \tag{MA-HM.15}
\]
Probability Haar on $C_K^1$ is understood.  Then
\[
       \operatorname{Tr}_{Z}C_g^Z
                 =\chi_{\pi_-}(f_{g,a}),
 \tag{MA-HM.16}
\]
where $\pi_-$ is Meyer's actual quotient
$H_-/(H_+\cap H_-)$, restricted to the indicated compact
character component.  In particular, this is an operator
trace comparison transported by MT-M24, not an inference
from equal spectra.

Let $\mathcal A_S(f)$ and $\mathcal D_S(f)$ denote the actual
traces in (MA-MEY.15):
\[
 \begin{aligned}
 \mathcal A_S(f)
   &=\operatorname{tr}_{\mathcal T}
           \lambda(f)(1-p_+^Si_+),\\
 \mathcal D_S(f)
   &=\operatorname{tr}_{\mathrm{common}}
           \lambda(f)(P_- -P_+^S).
 \end{aligned}
 \tag{MA-HM.17}
\]
The first is the local Fourier-cutoff trace and the second
is Meyer's common-space defect.  With all source domains
and nuclear extensions as in MA-MEY-003, the exact identity is
\[
       \boxed{\displaystyle
       \operatorname{Tr}_{Z}C_g^Z
          =\chi_{\pi_+}(f_{g,a})
                 -\mathcal A_S(f_{g,a})
                 +\mathcal D_S(f_{g,a}).}
 \tag{MA-HM.18}
\]
Over $\mathbb Q$, on the trivial compact character component
and with $a=1/2$, (MA-MEY.12) computes the defect as the
negative omitted-prime contribution.  Therefore
\[
       \operatorname{Tr}_{Z}C_g^Z
           =P(g)-W_S(f_{g,1/2})-B_{\overline S}(g).
 \tag{MA-HM.19}
\]
All three expressions on the right come from the actual
Meyer maps.  The omitted-prime term is the trace-class
Euler commutator of MA-MEY-001, with the sign fixed from
the definitions of $P_+^S$ and $P_-$.

\emph{Proof.}
Apply $q$ to the nuclear expansion in the proof of
MA-HM-001 and denote its bounded vectors in $Q$ by
$u_{n,j}=qv_{n,j}$.  The factorization gives
\[
 \begin{aligned}
 i_aS_g&=\sum_{n,j}\lambda_{n,j}
                     (i_au_{n,j})\otimes\ell_{n,j},\\
 S_gi_a&=\sum_{n,j}\lambda_{n,j}
                      u_{n,j}\otimes(\ell_{n,j}i_a).
 \end{aligned}
 \tag{MA-HM.20}
\]
The first series converges in Hilbert trace norm: $i_a$
maps the bounded vector family into a bounded subset of
$Z$, the functionals have uniformly bounded norm, and the
coefficients are in $\ell^1$.  The second is an order-zero
nuclear representation on $Q$, because the composed
functionals are equicontinuous and the target vectors are
bounded.  Each scalar trace series converges absolutely,
and both traces are
\[
          \sum_{n,j}\lambda_{n,j}
                            \ell_{n,j}(i_au_{n,j}).
 \tag{MA-HM.21}
\]
This proves cyclicity here directly from one common
rank-one representation.  It does not appeal to a trace
identity for arbitrary bounded lifts on $H$.

For completeness, the trace on $Q$ is unambiguous.
The order-zero framework is Meyer (2005), \S2.1, in
particular Theorem~2.1 and the trace pairing preceding it.
One can also see the needed approximation property
directly from MT-M18.  On $s(J)$ let $E_N$ keep the
blocks with $|n|\le N$.  For a bounded subset
$\mathcal U\subset s(J)$,
\[
    \sup_{z\in\mathcal U}q_k((1-E_N)z)
       \le (N+2)^{-1}\sup_{z\in\mathcal U}q_{k+1}(z)
                        \longrightarrow0.
 \tag{MA-HM.22}
\]
Thus finite-rank maps approximate the identity uniformly
on bounded sets.  The canonical nuclear trace agrees
with the trace obtained from these finite-rank
approximations.  Equations (MA-HM.6) and (MA-HM.21)
give (MA-HM.14).

The topological and bounded-bornological isomorphism in
MT-M22/MT-M24 sends $[F]\in Q$ to the class of
$b\mapsto\beta(b)^{-1}F(\log|b|)$ in the actual Meyer
negative quotient.  Its translation formula is
(M22.4).  Integrating that formula against (MA-HM.15)
cancels the character $\beta$ and leaves convolution by
$e^{-ax}g(x)=k_g(x)$.  The compact group has volume one.
Thus it carries $C_{k_g}^Q$ to the actual integrated
operator $\pi_-(f_{g,a})$.
It also carries the bounded vectors and equicontinuous
functionals in (MA-HM.20) to corresponding ones, preserving
each rank-one trace pairing.  This proves (MA-HM.16)
without using the zero-divisor trace formula.

Finally the actual trace additivity in the proof of
Meyer Theorem~5.8, author pp.~43--44, is exactly
\[
       \chi_{\pi_+}(f)-\chi_{\pi_-}(f)
                         =\mathcal A_S(f)-\mathcal D_S(f).
 \tag{MA-HM.23}
\]
Its spaces, maps and fixed-$S$ correction are recorded
and checked in MA-MEY.13--MA-MEY.16.  Rearrangement with
(MA-HM.16) proves (MA-HM.18), and (MA-MEY.12) proves
(MA-HM.19).  The polar trace for the trivial rational
character is
$P(g)=\int_{\mathbb R}g(x)(e^{-x/2}+e^{x/2})\,dx$.
For a nontrivial finite-order compact character the
polar quotient is killed by character projection.
\hfill$\square$

\subsection{A structural obstruction on a genuinely partial strip}

\paragraph{MA-HM-003 (exterior jet blocks are lost before any smoothing).}
Retain the spaces in (MA-HM.1), but do not assume (MA-HM.2).
The natural map $i_a:Q\to Z$ is still continuous with dense
image.  For each nontrivial completed zero $\rho$, let
$Q_\rho\subset Q$ be its finite-dimensional local jet block,
of dimension $m_\rho$, supplied by MT-M18 and the Chinese
remainder decomposition of its finite block algebra.
If
\[
              \operatorname{Re}(\rho-a)\notin[u,v],
 \tag{MA-HM.24}
\]
then
\[
                         i_a(Q_\rho)=0.
 \tag{MA-HM.25}
\]
This statement concerns zeros strictly outside the closed
strip.  It makes no assertion that boundary jets are lost.

Consequently, if a linear map $T_g:Z\to Q$, even without a
continuity assumption, is to satisfy
\[
                    T_gi_a=C_{e^{-a\cdot}g}^{Q},
 \tag{MA-HM.26}
\]
then at every zero in (MA-HM.24) it is necessary that
\[
               (\mathcal Mg)^{(j)}(\rho-a)=0
                       \quad(0\le j<m_\rho).
 \tag{MA-HM.27}
\]
For any one such zero there exists
$g\in C_c^\infty(\mathbb R)$, identically zero near $0$,
with $\mathcal Mg(\rho-a)=1$; for that allowed test
no map satisfying (MA-HM.26) exists.

In particular, for $a=1/2$ the following property is
equivalent to RH for this one function: for every
$\alpha>0$ and every compact smooth $g$ with $g(0)=0$,
the actual quotient $Z_{1/2,[-\alpha,\alpha],s}$ admits
the two smoothing factorizations (MA-HM.6) through
the full smooth quotient $Q$.  Under RH those maps
can be taken of order zero by MA-HM-001; the reverse
implication already follows from (MA-HM.26).

\emph{Proof.}
Continuity and density of $i_a$ use only the definitions:
$T_a:\mathcal E\to H$ is continuous, its image of
$\mathcal I_\eta$ lies in $J$, and its range contains
the compact smooth functions.  The ideal and its
Hilbert closure are invariant under all translations.
For the actual weight in (MA-HM.1) the sharper form
of the elementary translation estimate is
\[
 \|U_y\|_{H\to H}
 \le 2^{s/4}\max(e^{uy},e^{vy})(1+y^2)^{s/4}.
 \tag{MA-HM.28}
\]
Indeed the exponential part of $w(x+y)/w(x)$ is at
most $\max(e^{2uy},e^{2vy})$, and
$1+(x+y)^2\le2(1+x^2)(1+y^2)$ controls the polynomial
part.  The same bound holds for the induced group
on $Z$.  No zero-location assertion enters this bound.

On $Q_\rho$, $U_y^Q$ is multiplication by the class
of $e^{wy}$ in the local algebra at $w=\rho$.
After the intertwining identity
$U_y^Zi_a=e^{-ay}i_aU_y^Q$, its image
$Y_\rho=i_a(Q_\rho)$ is a finite-dimensional invariant
space.  If it is nonzero, the sole eigenvalue of
$U_1^Z$ on it is $e^{\rho-a}$: it is a nonzero
quotient of the generalized eigenspace with precisely
that eigenvalue.  Choose an eigenvector in $Y_\rho$.
For positive integers $n$, (MA-HM.28) then implies
\[
       e^{n\operatorname{Re}(\rho-a)}
          \le 2^{s/4}e^{vn}(1+n^2)^{s/4}.
 \tag{MA-HM.29}
\]
This is impossible when
$\operatorname{Re}(\rho-a)>v$.
Applying the same argument to $U_{-1}$ makes it
impossible when $\operatorname{Re}(\rho-a)<u$.
Hence $Y_\rho=0$, proving (MA-HM.25) for the entire
local jet algebra, not only its reduced eigenline.

Equation (MA-HM.26) would force
$C_{e^{-a\cdot}g}^Q$ to vanish on each such $Q_\rho$.
There it is multiplication by the class of
$\mathcal Mg(w-a)$ in
$\mathbb C[w]/(w-\rho)^{m_\rho}$.
Multiplication is the zero operator exactly when
this class is zero, which is (MA-HM.27).
Choose $b\in C_c^\infty((1,2))$ with $\int b=1$
and set $g(x)=e^{-(\rho-a)x}b(x)$.
Then $g$ vanishes near $0$ and
$\mathcal Mg(\rho-a)=1$, giving the stated obstruction.

If a zero is off the critical line, choose
$0<\alpha<|\operatorname{Re}\rho-1/2|$ and apply
this test to its exterior block.  Thus the
every-positive-width factorization property forces
RH.  Conversely RH supplies (MA-HM.2) with
$c_0=c_1=1/2$ for every such width, and MA-HM-001
constructs the required maps.
\hfill$\square$

The obstruction is loss of an invariant finite-length
spectral block under the actual completion map.
It is not an obstruction inferred from equal or
unequal scalar traces.  It also does not assert that
no useful map can factor through a quotient retaining
only the interior blocks.  Such a quotient has
discarded the exterior information in the full Meyer
quotient; composing it with smoothing or the Euler
commutator cannot recover that information by a
return map satisfying (MA-HM.26).
Boundary behavior remains separate, as in MA-TR-103.

\subsection{The actual cutoff limit after the stated strip premise}

\paragraph{MA-HM-004 (completed operator transport under the containing-strip premise).}
Specialize to the trivial rational character, $a=1/2$,
$I=[-\alpha,\alpha]$, and assume (MA-HM.2).
Use the literal orthogonal projections $P_Z$ and
$P_G=1-P_Z$ of the corrected Hilbert model.
Let $g\in C_c^\infty(\mathbb R)$ satisfy $g(0)=0$,
put $h(t)=g(\log t)$, and let
$\rho_\nu\in C_c^\infty(\mathbb R)$ be uniformly
bounded and tend to $1$ almost everywhere.
Then the actual geometric-order operators are
trace class and
\[
 \begin{aligned}
 \lim_\nu\operatorname{Tr}_{H}
                  (M_{\rho_\nu}C_gP_G)
   &=\mathcal A_S(f_{g,1/2})
          -\mathcal D_S(f_{g,1/2})
          -\chi_{\pi_+}(f_{g,1/2})\\
   &=W_S(f_{g,1/2})+B_{\overline S}(g)-P(g)\\
   &=A_0(h)+B(h)-P(h).
 \end{aligned}
 \tag{MA-HM.30}
\]
Every equality is obtained from traces of the actual
operators described in MA-LIFT-001--003,
MA-HM-001--002, and MA-MEY-001--003.
The proof does not substitute an equality of scalar
zero sums.  It gives a trace-limit comparison;
it does not identify each individual finite
orthogonal cutoff operator with a local Fourier
cutoff operator.

\emph{Proof.}
MA-LIFT-002 proves that the literal ambient operator
$C_gP_Z$ is trace class and that its trace equals
the quotient convolution trace.  We use that
operator assertion and its proof by trace cyclicity,
not the subsequent evaluation of the quotient
trace as a zero sum.  MA-LIFT-003 proves
\[
 \|(M_{\rho_\nu}-1)C_gP_Z\|_1\longrightarrow0,
 \qquad
 \operatorname{Tr}_{H}(M_\rho C_g)
                         =g(0)\int_{\mathbb R}\rho(x)\,dx.
 \tag{MA-HM.31}
\]
Thus, since $g(0)=0$,
\[
 \lim_\nu\operatorname{Tr}_{H}
                  (M_{\rho_\nu}C_gP_G)
       =-\operatorname{Tr}_{H}(C_gP_Z)
       =-\operatorname{Tr}_{Z}(C_g^Z).
 \tag{MA-HM.32}
\]
These are limits of traces of the actual trace-class
operators in the order specified by the draft.
Only the spectral product in (MA-HM.31) is asserted
to converge in trace norm.  The cusp of the
weight and its finite boundary contribution were
included in MA-LIFT-002.

Now apply the operator trace comparison
(MA-HM.18).  Its negative is the first line of
(MA-HM.30).  The common-space defect is
$\mathcal D_S=-B_{\overline S}$ by MA-MEY-003,
with the omitted-prime contribution computed
from the trace-class commutator in MA-MEY-001.
The polar quotient has trace $P(g)$.
Finally the independently computed local
Fourier-cutoff trace is
$W_S=A_0(h)+B_S(h)$ when $h(1)=0$.
Combining its finite prime terms with the
omitted-prime trace gives $B(h)$, proving the
last line.  This local evaluation is the one
used in MA-MEY.15--MA-MEY.16; it requires no
zero-location input.
\hfill$\square$

For every rational or Dirichlet narrow strip,
the operator part of this transport can therefore
be instantiated using $A$.  On unconditional
wide strips it needs no such input.  The
operator bridge is complete under its stated
premise; obtaining that premise on arbitrarily
narrow strips without using zero location is
the remaining independent-method question.

\subsection{What this bridge settles and what it does not}

The smoothed Hilbert operator now has an explicit route to
the genuine common-space geometry:
$C_g^Z=i_aS_g$, $S_gi_a=C_{k_g}^Q$, the MT-M24 map identifies
that quotient operator with $\pi_-(f_{g,a})$, and
(MA-HM.18) expresses its trace by the actual local and
arithmetic-defect operators.  The route carries prime
contributions through MA-MEY-001/003 and makes no demand
that a narrow-strip M\"obius inverse be bounded.

The only zero-location premise in this route is
(MA-HM.2), used in obtaining the uniformly bounded
Hilbert restrictions and the interior adapted blocks.
It is harmless on wide strips.  Using $A$ to verify it
on every Dirichlet narrow strip is a valid conditional
application of the already granted theorem, but is
circular in a proposed zero-location-free proof of the
method under test.  For the imaginary quadratic
character of MA-CH-003, the corresponding invocation
needs its own RH, sufficiently $M$.

MA-LIFT-001--003 and MA-HM-004 also give the
specific ambient cutoff trace limit under that same
premise.  This discharges the ambient comparison
for the corrected model at the level asserted in
(MA-HM.30).  No equality of individual finite
cutoffs is inferred.  In particular, a compact
cutoff cannot silently replace Meyer's one-sided
commutator: MA-MEY-002 gives an explicit failure
of trace-norm convergence for that substitution.
MA-HM-003 explains structurally why a completion
that genuinely drops exterior zero blocks cannot
be used to recover the full smooth quotient by
the same smoothing factorization.

The functional-analytic location is Grothendieck's
nuclear tensor product and trace pairing: the common
rank-one expansion (MA-HM.20) is the concrete object
that makes the trace comparison valid.  It fits
the inherited cote~1 determinant/trace programme,
scans~079--080 and~111--112, through ALPOGE\_BRIDGES
BR-TR-01/02 and F-SM-11/12.  Those archive placements
are carried from the supplied runs; no new archive
scan reading is claimed here.  The quotient still
records the full finite-length spectral jet algebras.
Neither this nuclear comparison nor its Hilbert
factorization supplies faithful positivity, replaces
those jets by a reduced divisor, or realizes them
as the cohomology of a proper scheme.

\paragraph{Exact source scope.}
The inherited analytic data are GRH\_NIGHT MT-M18
(M18.1)--(M18.4), (M18.13), and MT-M26
(M26.2), (M26.5)--(M26.12); the actual smooth-quotient
comparison is MT-M22 (M22.2)--(M22.4), (M22.10) and
MT-M24 (M24.3), (M24.6)--(M24.9).
Meyer's published source is
\emph{On a representation of the idele class group related
to primes and zeros of $L$-functions},
\emph{Duke Mathematical Journal, 127} (2005), 519--595,
arXiv:math/0311468v3: \S2.1 for the nuclear trace
pairing and Theorem~2.1; Lemmas~5.4--5.6 for the
common-space maps and nuclear defect; Theorem~5.8,
author pp.~43--44, for trace additivity.
Its Theorem~5.11 or the resulting scalar zero sum
is not needed to establish the operator trace
identity (MA-HM.18).
